\section{Introduction}
\label{sec:intro}

Trillium's B200 partition gives exclusive access to 52 nodes of 8 B200 GPUs for a few days
\est{}. We use the window to measure how well the partition trains large language models, and which
configurations make the best use of it. \Cref{tab:machine} lists what is known about the machine before
access.

\begin{table}[h]
  \caption{Trillium's B200 partition, as known before access.}
  \label{tab:machine}
  \centering\small
  \begin{adjustbox}{max width=\linewidth}
  \begin{tabular}{lll}
    \toprule
    & Known & Source \\
    \midrule
    Nodes & 52 (\texttt{trig0101}--\texttt{trig0152}), exclusive access during the window & SciNet talk 2026-10-21; \texttt{/etc/hosts} \\
    Per node & 8 B200 180\,GB on NVLink, 2 $\times$ 48-core Intel CPUs, 2\,TB RAM, liquid cooled & SciNet talk \\
    GPUs & 416 in total & \\
    Network & InfiniBand, 800\,Gb/s per node, ``non-blocking''; number of HCAs unknown & user; SciNet \\
    Switch layout & unknown; \texttt{/etc/hosts} puts the nodes in 8 subnets of 7, 7, 7, 7, 4, 7, 7, 6 nodes & \texttt{/etc/hosts} \\
    Shared with others & distributed filesystem and interconnect (CPU and H100 nodes stay in production) & user \\
    Software & \texttt{module load cuda/12.9}, \texttt{--gpus-per-node=b200:8} & SciNet talk \\
    Matmul ($8192^2$) & B200: BF16 1{,}668, FP8 3{,}041\,TFLOP/s; H100: BF16 785\,TFLOP/s & SciNet (their measurement) \\
    Window & a few days \est{}, date unknown & user \\
    \bottomrule
  \end{tabular}
  \end{adjustbox}
\end{table}

\paragraph{What sets the partition apart.}
\begin{enumerate}[leftmargin=*, itemsep=1pt]
  \item \textbf{Network per FLOP.} Each GPU has 100\,Gb/s of inter-node bandwidth for 1{,}668\,TFLOP/s
        of BF16 matmul (SciNet's measurement): about 0.060\,Gb/s per TFLOP/s. This is close to Rorqual
        (50\,Gb/s for 785\,TFLOP/s, 0.064), where the network limited training, but the NVLink domain is
        twice as large (8 GPUs instead of 4). A standard DGX B200 has 4 times more network per node.
        Configurations that keep most traffic inside a node should do best.
  \item \textbf{FP8 and MXFP8.} Low precision doubles the matmul rate (3{,}041\,TFLOP/s in FP8, SciNet's
        measurement) but not the network. What is network-bound in BF16 becomes more so.
  \item \textbf{Memory.} 180\,GB per GPU, 73\,TB in total. Qwen3.8-27B (about 430\,GB of training state
        \est{}) fits on one node; GLM-5 fits on about 13 nodes \est{}, so up to 4 copies run at once.
        Activation checkpointing can be reduced.
  \item \textbf{Unknown topology, shared fabric.} The switch layout is not known, and the filesystem
        and interconnect carry other users' traffic. Placement enters the design only if the pairwise
        measurements show a structure (\cref{sec:network}); a reference run repeated through the window
        separates our configurations from background load (\cref{sec:interference}).
\end{enumerate}

\paragraph{Questions.} Each section answers one question:
(\cref{sec:network}) are all GPUs and links alike, and does the fabric have a structure?
(\cref{sec:interference}) can jobs run side by side without slowing each other?
(\cref{sec:layouts}) which parallel layout should each workload use?
(\cref{sec:scaling}) how far does each workload scale on the partition?
(\cref{sec:precision}) what do FP8 and MXFP8 buy, and do they train the same?
(\cref{sec:operations}) what do checkpoints, startup and energy cost, and is throughput stable?

\paragraph{Declared deviations} from the reference training recipes, valid for every result unless noted:
\placeholder{list once the B200 configurations are fixed (optimizer, multi-token prediction, data).}
